\labsection{1}{Frame a question and build a reproducible data profile}{sec:lab01}

\begin{labproject}{Lab 1 guided workflow}
\textbf{Dataset.} \texttt{bike\_sharing\_dataset\_day.csv}.\\
\textbf{Goal.} Convert a real operational dataset into a traceable question--evidence workflow.

The daily bike-sharing table contains one row per calendar day, weather/context variables, and the total rental count \rcode{cnt}. It is a useful first dataset because students must recognise \textbf{time}, \textbf{measurement grain}, and \textbf{context} before calculating anything.

\textbf{Procedure.}
\begin{enumerate}
  \item State the unit of observation and distinguish identifiers, contextual variables, predictors, and outcomes.
  \item Parse \rcode{dteday} as a date and inspect dimensions, column types, ranges, and missingness.
  \item Write one descriptive, one comparative, and one relational question before computing any result.
  \item Compare rental demand between working and non-working days and examine the relationship between temperature and demand.
  \item Plot demand over time and identify at least two features that a single overall mean would hide.
  \item Record \rcode{sessionInfo()} so the computational environment becomes part of the evidence record.
\end{enumerate}

\textbf{Interpretation checkpoint.} A correlation between temperature and rentals does not by itself show that changing temperature causes demand to change; season, month, weather state, and long-term trend are competing explanations.

\medskip
\textbf{Executable experiment.} Run and edit the file below from a clean R session.
\lstinputlisting[style=dsr]{all_experiments/lab01_workflow.R}
\end{labproject}
